\begin{definition}[Marked surfaces, local covers, localization; after Kahn]\label{def:localization-seed}
A \emph{marked Riemann surface} is a pair \((S,A)\) with \(S\) a Riemann surface and \(A\subset S\) compact, connected, and full in the sense that no component of \(S\setminus A\) has compact closure in \(S\); a morphism \(g\colon(S,A)\to(T,B)\) is a holomorphic map with \(A=g^{-1}(B)\).  The pair is \emph{local} if every component of \(S\setminus A\) is an annulus; a \emph{local neighborhood} of \(A\) is a subsurface \(U\supset A\) with \((U,A)\) local.  A morphism is a \emph{local cover} if it restricts to a bijection of the marked continua and to a covering map of their complements.  Kahn proves that every \((S,A)\) admits a \emph{localization}: a local pair \((S_A,\widehat A)\) with a local cover \(\pi_{S,A}\colon(S_A,\widehat A)\to(S,A)\), unique up to isomorphism.  When \(S\setminus A\) is connected, \(S_A\setminus\widehat A\) is the cover of \(S\setminus A\) corresponding to the image of \(\pi_1(U\setminus A)\) for a local neighborhood \(U\) of \(A\), completed by regluing \(U\) along its natural lift.  Every morphism \(g\) has a unique commuting localized lift \(\widehat g\), and \(\widehat g\) is an immersion whenever \(g\) is (Kahn~I, \S7.4.1, Theorem~7.12, pp.~43--45, arXiv:math/0609045v3).  The source asserts preservation of \emph{immersion}; it never asserts preservation of embedding.
\end{definition}

\begin{definition}[Dynamical localization and canonical renormalization]\label{def:canonical}
Given a \(\psi\)-polynomial-like \((i,f)\colon(\mathbf U',K')\to(\mathbf U,K)\) of degree \(d\), an FJ-set \(A\subset K\) with \(A'=i^{-1}(A)\),\footnote{At Kahn~I, \S7.4.2, p.~45 the hypothesis list for dynamical localization reads \(A'=i^{-1}(A)\), and \S7.4.3, p.~45 then writes \(K_p'=i^{-1}(K_p)\) while simultaneously instructing the reader to replace \((i,f)\) by the iterate \((i^p,f^p)\colon\mathbf U^p\to\mathbf U^0\).  Read with the un-iterated \(i\), the two prescriptions name different maps.  We use throughout the intended reading, in which \(i\) denotes the inclusion leg of the map actually being localized: \(A'=(i^p)^{-1}(A)=A\) at the ordinary primitive seed, since \(K_p\subset U^p\).  This matters because Definition~\ref{def:excess} turns on distinguishing \(A'\) from the other components \(X\) of \(F^{-1}(A)\).} and a compact \(B\subset K\) satisfying Kahn's hypotheses (in particular \(A\) is a component of \(B\), \(B\) is compact with \(\mathbf U\setminus B\) connected, \(A'\) is a component of \(f^{-1}(A)\), and all branch values of \(f\) lie in \(B\)), set
\[
 T=(\mathbf U\setminus B)\cup A,\qquad T'=(\mathbf U'\setminus f^{-1}B)\cup A'.
\]
Localizing the restricted morphisms yields the \emph{dynamical localization} \((\widehat i,\widehat f)\colon(T'_{A'},\widehat{A'})\to(T_A,\widehat A)\), a \(\psi\)-polynomial-like map in which \(\widehat i\) is an immersion and \(\widehat f\) is a branched cover of degree equal to the local degree of \(f\) at \(A'\) (Kahn~I, \S7.4.2, p.~45).  The \emph{canonical renormalization} of a primitively \(p\)-renormalizable \(\psi\)-quadratic-like map is the dynamical localization of \((i^p,f^p)\) at \(A=K_p\) with removal of \(B=\mathcal K_p\); its central local degree is two (Kahn~I, \S7.4.3, p.~45).  For an ordinary primitive seed the marked surfaces are exactly \(T'\), \(T\) of (1.1), with \(A'=A\).

This choice of \(B\) is not the only one available in the source, and the distinction is load-bearing here.  Kahn~I, \S7.4.4, p.~46 describes two \emph{further} natural renormalizations obtained as dynamical localizations of the \emph{same} iterate \((i^p,f^p)\) at the \emph{same} \(A=K_p\), but with a different removal set: one takes \(B:=A\cup\mathcal K_{pq}\), and a second variant is built from the postcritical set and the geodesic \(\gamma_p\) (it is for that variant that the \(pq\)-renormalization is the \(q\)-renormalization of the \(p\)-renormalization).  Throughout this paper ``canonical renormalization'' always means the \S7.4.3 construction with \(B=\mathcal K_p\).  The restriction is essential rather than decorative: the excess-component and removal-asymmetry argument of Section~\ref{sec:forced} turns on the source removing all of \(F^{-1}B\) while the target removes only \(B\), and enlarging \(B\) to \(A\cup\mathcal K_{pq}\) changes that asymmetry.  We make no claim about the \S7.4.4 variants.
\end{definition}

\begin{definition}[Excess component]\label{def:excess}
An \emph{excess component} is a component \(X\ne A\) of \(F^{-1}A\).  (For the physical inclusion \(A'=(i^p)^{-1}A=A\); the selected \(X\) is a \emph{different} component of \(F^{-1}A\) and must not be confused with \(A'\).)
\end{definition}

\begin{definition}[Direct-decoration alignment]\label{def:decoration}
A canonical seed \emph{carries KL direct-decoration data} with \emph{decoration parameters} \((q,\nu)\) --- KL's own term (KL~II, Corollary~5.14 and Theorem~5.15, pp.~23--24) --- if the Yoccoz-puzzle apparatus of KL~II applies to it, producing \(P=Y^{(\nu-1)q+1}\) and the two univalent \(f^q\)-pullbacks \(Q_L,Q_R\subset P\) of depth \(\nu q+1\) (KL~II, \S2.1, p.~6, and Figure~2.1, p.~7).  The \emph{exact direct-decoration depths} are \(m=\nu q+1\) and \(n=(\nu-1)q+1\).
\end{definition}